\chapter{Introduction}
\label{chap:ch1}

The ability to look at a collection of broken fragments and reason about how they once fit together is a form of spatial intelligence that humans have practiced for millennia. In archaeology, researchers reassemble eroded pottery, shattered frescoes, and fragmented statues to preserve the material record of ancient civilizations \cite{huang2006reassembling, toler2010multi, funkhouser2011learning}. In paleontology, fractured bones and lithic artifacts are pieced together to reconstruct tools and skeletal remains. In medicine, surgeons mentally reconstruct comminuted bone fractures before planning an operation.

Despite the importance of these tasks, they remain overwhelmingly manual. An experimental study on lithic refitting reported a success rate of only 30\%, with experts performing only marginally better than novices \cite{laughlin2010lithic}. Museum storerooms around the world hold thousands of unassembled fragments, waiting to be pieced together by specialists whose time is limited and whose accuracy, even with expertise, is modest. The sheer scale of the problem, combined with the risk of further damaging fragile artifacts through physical handling, has long motivated the search for computational alternatives.

With the advent of high-resolution 3D scanning, the physical pieces of a broken object can now be digitized into point clouds: unordered sets of 3D coordinates that capture the geometry of each fragment. Given a set of such point clouds, the goal is to predict the rigid transformation (a rotation and a translation) that maps each fragment back to its original position. The problem is known as 6-Degree-of-Freedom (6-DoF) pose estimation and sits at the intersection of computer vision, computational geometry, and deep learning.

Automating 3D reassembly is far from trivial. Unlike 2D jigsaw puzzles, where flat edges and printed patterns constrain the solution, 3D fracture fragments feature irregular, continuous surfaces with no discrete boundary markers. The search space grows combinatorially with the number of pieces, and small errors in local pairwise alignments propagate and accumulate when assembling multi-part objects, a phenomenon known as drift.

Early computational approaches relied on hand-crafted geometric descriptors (integral invariants \cite{huang2006reassembling}, curvature profiles, surface normals, and boundary contours \cite{toler2010multi}) to find matches between fracture surfaces. While effective on clean, controlled breaks, these methods are sensitive to surface ambiguity and deteriorate when applied to eroded, incomplete, or noisy real-world fragments. A parallel line of research addressed semantic part assembly, reconstructing objects from predefined, functionally meaningful components such as furniture legs or chair backs \cite{Huang2020DGL, harish2022rgl, wu2020pq}. These methods leverage categorical priors and part labels, but their assumptions break down when applied to arbitrary fracture fragments that carry no semantic identity.

The landscape shifted with two key developments. First, specialized neural architectures for unstructured 3D data, PointNet \cite{qi2017pointnet} and PointNet++ \cite{qi2017pointnetplusplus}, enabled the direct processing of point clouds, replacing hand-crafted features with learned geometric representations. Second, large-scale fracture datasets, most notably Breaking Bad \cite{sellan2022breaking}, provided the training data necessary for data-driven approaches to geometric assembly. Building on these foundations, the Jigsaw framework \cite{lu2023jigsaw} established the first learning-based pipeline for multi-piece fracture assembly, combining hierarchical point cloud features with attention mechanisms \cite{vaswani2017attention}, global bipartite matching via Sinkhorn optimal transport \cite{cuturi2013sinkhorn}, and transformation synchronization through Shonan averaging \cite{dellaert2020shonan}. Jigsaw set the standard baseline for subsequent research on Breaking Bad.

Since Jigsaw, the field has evolved rapidly along several directions. To address the computational cost of dense point matching, PMTR \cite{lee20243d} introduced sub-quadratic proxy tensors, while PHFormer \cite{cui2024phformer} decomposed pose estimation into adjacency prediction and transformation regression using proxy-level attention. To enrich attention with spatial reasoning, the Geometric Point Attention Transformer (GPAT) \cite{Li2023GPAT} demonstrated that encoding pairwise distances and angles between fragments directly into the attention scores can improve assembly quality. To introduce global shape awareness, Jigsaw++ \cite{lu2024jigsawpp} incorporated a shape prior through latent space reconstruction using Rectified Flow. To bridge the gap between synthetic training data and real-world fractures, GARF \cite{Li2025GARF} recast the problem as learning continuous velocity fields on the Lie group $SE(3)$, achieving state-of-the-art performance by moving fragments along geodesic paths.

Despite this rapid progress, several open questions remain within the Jigsaw pipeline, which continues to serve as the foundational baseline for fracture assembly on Breaking Bad. Its cross-attention mechanism can produce overconfident correspondence distributions, particularly on out-of-distribution objects, and may benefit from regularization techniques that smooth the attention weights. The feature extractor uses a single round of attention, while recent work (PHFormer, GARF) suggests that iterative refinement of self-attention and cross-attention improves geometric representations. The PointNet++ backbone constructs local neighborhoods via $k$-nearest neighbors without regard for surface topology, potentially including spatially close but topologically distant points in each centroid's neighborhood. And pairwise pose estimates are recovered from sparse RANSAC correspondences without fully exploiting the dense geometry available on fracture surfaces.

\section{Objectives}

The objective of this thesis is to investigate whether targeted, modular modifications to the Jigsaw framework can address the limitations outlined above, and what such experiments reveal about the relationship between architecture, data, and optimization in 3D fracture assembly. Rather than proposing a new architecture, we study four specific extensions: a geometric attention bias for cross-attention, a gated second attention pass, a Gabriel graph neighborhood filter for PointNet++, and an ICP refinement step for pose recovery; alongside a data scaling analysis and an interactive web application for visual comparison of results.

\section{Original Contributions}

The main contributions of this thesis are:

\begin{enumerate}
    \item \textbf{Geometric pair attention as learned soft regularization.} We integrate the pair attention mechanism from GPAT into Jigsaw's cross-attention layers, where it learns a data-dependent additive bias that softens the attention distribution. This acts as a learned regularizer that consistently improves out-of-distribution generalization on the artifact subset while remaining neutral on in-distribution performance, a pattern that holds across data regimes and piece counts.

    \item \textbf{Gated double attention with training dynamics analysis.} We extend the feature extractor with a second attention pass using zero-initialized ReZero gates. On a reduced dataset the additional capacity acts as an implicit regularizer, providing +2 to 5 percentage points. On the full dataset, the baseline already captures these patterns, confirming that data and architecture address the same representational bottleneck from complementary directions. We further analyze the learned gate values and demonstrate the sensitivity of gated finetuning to loss scheduling.

    \item \textbf{Gabriel graph neighborhood filtering.} We replace the standard kNN grouping in PointNet++ with a Gabriel graph filter that prunes topologically distant neighbors from each centroid's neighborhood. The experiment shows that surface-faithful neighborhoods produce higher-quality matches and improve cross-domain pose metrics on artifact objects, but the reduced neighborhood diversity slightly lowers the overall correspondence count. A systematic study of the retention ratio $r$ reveals the trade-off between pruning aggressiveness and feature aggregation quality.

    \item \textbf{ICP refinement for pairwise pose estimation.} Point-to-point ICP inserted between RANSAC and Shonan averaging yields a consistent improvement of approximately 1 percentage point across all configurations. The modification is non-learned, orthogonal to the attention-based extensions, and carries no risk of overfitting.

    \item \textbf{Data scaling analysis.} By comparing results on the full dataset against a reduced subset ($\sim$63\% of everyday training patterns), we quantify the effect of data quantity on assembly quality (the full dataset yields over 10 percentage points more) and show that the proposed architectural modifications provide consistent improvements under data scarcity, acting as implicit regularizers that complement the baseline.

    \item \textbf{Interactive web application.} A Django and FastAPI-based demo with Three.js 3D visualization, enabling users to run inference across eleven model variants and visually compare reconstruction quality.
\end{enumerate}

\section{Thesis Structure}

The remainder of this thesis is organized as follows:

\begin{itemize}
    \item \textbf{Chapter~\ref{chap:ch2}} defines the fracture assembly problem.

    \item \textbf{Chapter~\ref{chap:ch3}} covers the theoretical foundations (point cloud processing, attention mechanisms, pose estimation).

    \item \textbf{Chapter~\ref{chap:ch4}} reviews related work.

    \item \textbf{Chapter~\ref{chap:ch5}} presents the proposed extensions to the Jigsaw pipeline.

    \item \textbf{Chapter~\ref{chap:ch6}} reports the experimental evaluation.

    \item \textbf{Chapter~\ref{chap:ch7}} describes the interactive web application.

    \item \textbf{Chapter~\ref{chap:ch8}} concludes with findings, threats to validity, and future work.

\end{itemize}
